\documentclass[11pt]{article}
\usepackage[margin=1.05in]{geometry}
\usepackage{amsmath,amsthm,amssymb,mathtools,bm}
\usepackage{booktabs,array,enumitem}\usepackage[expansion=false]{microtype}
\usepackage{xcolor}
\usepackage[colorlinks=true,linkcolor=blue!60!black,citecolor=blue!60!black,urlcolor=blue!60!black]{hyperref}
\theoremstyle{plain}
\newtheorem{theorem}{Theorem}[section]\newtheorem{proposition}[theorem]{Proposition}\newtheorem{lemma}[theorem]{Lemma}
\newtheorem{corollary}[theorem]{Corollary}\newtheorem{conjecture}[theorem]{Conjecture}
\theoremstyle{definition}\newtheorem{definition}[theorem]{Definition}\newtheorem{remark}[theorem]{Remark}
\newtheorem{example}[theorem]{Example}
\newcommand{\R}{\mathbb R}\newcommand{\Z}{\mathbb Z}\newcommand{\E}{\mathbb E}\newcommand{\Var}{\operatorname{Var}}\newcommand{\Cov}{\operatorname{Cov}}
\newcommand{\relu}{\operatorname{relu}}\newcommand{\diag}{\operatorname{diag}}\newcommand{\tr}{\operatorname{tr}}
\newcommand{\Tr}{\operatorname{Tr}}\newcommand{\1}{\mathbf 1}\newcommand{\cA}{\mathcal A}\newcommand{\cB}{\mathcal B}
\newcommand{\cF}{\mathcal F}\newcommand{\cT}{\mathcal T}\newcommand{\cH}{\mathcal H}\newcommand{\aff}{\mathfrak{aff}}
\newcommand{\He}{\mathrm{He}}\newcommand{\Fa}{\Phi_{\mathrm F}}\newcommand{\cZ}{\mathcal Z}
\title{\textsc{Localization in the commutant, stage 12:}\\ \large the network as its own two-tier system.\\
Self-localization and the internal Eldan clock, a hyperbolic upper tier, the gate gas,\\ and the memory as a cumulant series
stratified by border codimension}
\author{}
\date{10 October 2026}
\begin{document}
\maketitle
\begin{abstract}
Stages 9--11 studied the bias-free He-initialised \textsc{ReLU} network with Gaussian input inside a two-tier framework imposed from
outside. The lower tier is the network's gates; the upper tier is Eldan's stochastic localization of the input. Here we show that
the network carries the same two-tier structure \emph{internally}, and that the external framework is the internal one shifted in
time. (i)~\emph{Self-localization.} Each layer of a fresh network acts as a localization step. Through every row of $W_{l+1}$, the law
of $h_l$ looks like a Gaussian ball whose standardized centre is $N(0,t_l)$-distributed across rows, exactly as Eldan's localized
input looks to a first layer at time $t_l$. The internal clock is $t_l=|\E h_l|^2/\Tr\Cov h_l=\cos\theta_l/(1-\cos\theta_l)$, where
$\theta_l$ is the folding angle of independent inputs. It obeys a closed recursion, and in the Fatou chart of stage 11 it advances by
exactly one unit per layer. Depth and localization time are therefore one coordinate, Szekeres' regular iteration defines fractional
depth, and the stage-11 collapse theorem is this statement seen from outside. Measured on a width-1024 network, the field law,
gate variance, clock and log-norm variance match the predictions. (ii)~\emph{Cooling.} A localized state is the Gibbs state of
$H=|x|^2/2$ at $\beta=1+t$, so the network anneals itself on the schedule $T_l=1-\cos\theta_l\approx\theta_l^2/2$. At finite width
the temperature splits into an angular part and a radial part, $T_l\approx(1-\cos\theta_l)+V_l/4$, where $V_l$ is the across-input
variance of $\log|h_l|^2$. This variance saturates at $44.6/n$ because its increments decay like $\theta^2$. The radial part takes
over at Fatou time $\tau_*\approx2\sqrt n$. (iii)~\emph{A hyperbolic upper tier.} Under the Fisher--Rao metric, the per-neuron
Gaussian laws form a hyperbolic plane. Positive homogeneity is translation along the kink geodesic $\{\mu=0\}$, and rectification is
a skew product over the signed distance to it. Criticality means zero drift along this geodesic, and the collective (norm) mode is
the finite-width diffusion along it. Measured, this diffusion, not the birth of cumulants at the walls, is what makes deep
pre-activations non-Gaussian at width 1024: each neuron's skewness is linear in its field ($R^2\approx0.9$) and its excess kurtosis is
flat, both proportional to $V$ with one efficiency $\approx0.7$ at every depth. (iv)~\emph{The gate gas.} The lower tier is a dichotomized-Gaussian lattice gas. Its Ising
couplings are $c_{ab}\varphi(r_a)\varphi(r_b)/(p_aq_ap_bq_b)$, with a Mattis component whose stored pattern is the field itself.
Every wall-localized one-body quantity obeys a \emph{wall-density law}: it equals $(\int Q)\sin(\theta/2)/\sqrt\pi$, so it decays
harmonically in Fatou time and its total over depth grows like $\log\tau$. (v)~\emph{Tier coupling.} At every layer an exact Euler
split writes the mean as an upper-tier term plus a gate--field covariance. The non-Gaussian corrections form an Edgeworth series. We
prove that each Edgeworth term is a sum over strata of the wall arrangement, and that the joint cumulant entering a stratum is that
of exactly the fields whose walls meet there: the Kikuchi level equals the codimension. Normal crossings carry nothing. Two-body
memory lives only on the bent crossings, where stage 11 found the non-abelian isotropy $\aff(1)$. (vi)~\emph{Modular structure.}
The layer's Gaussian state is a KMS state on the fibre relation of $h_l$. Its extremal decomposition is the law of $h_l$, its
conditional expectation onto the centre is the posterior, and the centres form a decreasing tower. The forward state localizes
while the backward (sensitivity) state delocalizes with the power law $(\tau_l/\tau_L)^3$. Every claim is labelled as proved, proved
at infinite width, conjectural or an analogy. The identities are checked by \texttt{scripts/verify\_stage12.py}.
\end{abstract}
\tableofcontents
\section{The realisation in one page}
\label{sec:summary}
\paragraph{1. One layer is a two-tier system (Section~\ref{sec:cycle}).}
\begin{itemize}[leftmargin=*]
\item The upper tier is the quasi-free projection $(\mu_l,C_l)$ of the pre-activations.
\item The lower tier is the gate gas.
\item The up map is the gating $h=Dz$ followed by $W_{l+1}$, and the closure is the second-order moment chain.
\item The memory is the non-quasi-free part of the law.
\item An exact Euler split, $\E f=\langle\E J_l,\mu_l\rangle+\sum_a\Cov(J_{l,a},z_{l,a})$, separates the tiers at every layer. At the input the whole
mean is the lower-tier border functional of stage 11.
\end{itemize}

\paragraph{2. The network localizes itself (Section~\ref{sec:clock}).} Through every row of the next layer, the law of $h_l$ is a Gaussian
ball whose standardized centre is $N(0,t_l)$ across rows, with $t_l=|m_l|^2/\Tr\Sigma_l\to\cos\theta_l/(1-\cos\theta_l)$. This is exactly what a
first layer sees of Eldan's localized input at time $t_l$, with rows playing the role of localization paths. The clock obeys
$t'=\E m_1^2/\E s_1^2$, which is the folding map in disguise, and in the Fatou chart it advances by one per layer. Depth is localization
time: Szekeres' regular iteration gives fractional depth, and external localization to time $t$ is fractional depth
$\Fa(\arccos\frac t{1+t})-\Fa(\pi/2)$. Localized states are Gibbs states of $|x|^2/2$ at $\beta=1+t$, so the network anneals itself on
$T_l=1-\cos\theta_l\approx9\pi^2/2\tau_l^2$. At finite width a radial temperature $V_l/4$ (the log-normal norm) is added. It explains the
measured clock at $l=16$ ($0.0779$ against $0.0705+0.0080$) and takes over at $\tau_*\approx2\sqrt n$.

\paragraph{3. The upper tier is hyperbolic (Section~\ref{sec:hyper}).} With the Fisher--Rao metric the per-neuron laws form a hyperbolic
plane. The field is the signed distance to the kink geodesic $\{\mu=0\}$, $r=\sqrt2\sinh(d/\sqrt2)$. Positive homogeneity is translation
along that geodesic, and rectification is a skew product with base map $R=m_1/s_1>0$, which pushes every neuron to the positive side.
Criticality means zero drift along the geodesic. The finite-width diffusion along it has across-input increments
$\frac4n(\frac32-\frac{J_2(\theta)}{2\pi})\approx4\theta^2/n$, which are summable, so $nV_l\to44.6$ rather than growing like $5l$. The
collective mode (trace share $\approx tV/4$) is this diffusion. It compresses the field law and skews the pre-activations: skewness is
linear in the field ($R^2\approx0.9$) and the excess kurtosis is flat, both proportional to $V$.

\paragraph{4. The lower tier is a gate gas with a wall-density law (Section~\ref{sec:gates}).} The gas is a dichotomized Gaussian with Ising
couplings $c_{ab}\varphi_a\varphi_b/(p_aq_ap_bq_b)$, and its Mattis component stores the field as its pattern. Every wall-localized
one-body quantity equals $(\int Q)\sin(\theta_l/2)/\sqrt\pi\approx2.66(\int Q)/\tau_l$. The lower tier is therefore active
harmonically in Fatou time, and its accumulated activity grows like $\log\tau$. Hermite cancellation suppresses neuron averages of
the marginal memory, but not its mean square.

\paragraph{5. The memory is a cumulant series stratified by codimension (Section~\ref{sec:coupling}).} Each Edgeworth term is a sum over
chain strata of the wall arrangement, and the cumulant on a stratum is the joint cumulant of the fields whose walls meet there.
One-body cumulants live on walls and two-body cumulants on bent crossings: the Kikuchi level is the codimension. Normal crossings carry
nothing, so the Gaussian closure is exact for them, and pair memory sits exactly at stage 11's $\aff(1)$ points. The third-order
formula is checked to one standard error, and the strata cancel. Third cumulants are born at rate $0.718/\tau$.

\paragraph{6. Modular structure (Section~\ref{sec:modular}).} The Gaussian state is KMS on the fibre relation of $h_l$. Its extremal
decomposition is the law of the representation, the conditional expectation onto the centre is the posterior, and the centres
$L^\infty(h_l)$ form a decreasing tower. Coarse-graining the centre and cooling are interchangeable on the next layer. The forward state
localizes while the backward state delocalizes, as $\prod(1-\theta_{k-1}/\pi)\approx(\tau_l/\tau_L)^3$; this is confirmed on the width-1024
network near the output and is about $20\%$ faster near the input.

\section{Two tiers inside one layer}
\label{sec:cycle}
\paragraph{Notation.} As in stage 11: $W_l\in\R^{n\times n}$ with i.i.d.\ $N(0,2/n)$ entries, $h_0=x\sim\gamma=N(0,I_n)$, $z_l=W_lh_{l-1}$,
$h_l=\relu(z_l)$, gates $g_l=\1\{z_l>0\}$, $D_l=\diag(g_l)$. For a fixed network, $m_l=\E h_l$, $\Sigma_l=\Cov h_l$, $\mu_l=\E z_l=W_lm_{l-1}$
and $C_l=\Cov z_l=W_l\Sigma_{l-1}W_l^{\top}$; expectations are over $x$ unless stated. Neuron $a$ of layer $l$ has \emph{scale}
$s_a=C_{l,aa}^{1/2}$, \emph{field} $r_a=\mu_{l,a}/s_a$ and gate probability $p_a=\Pr(g_{l,a}=1)$, with $q_a=1-p_a$. The folding angles
are $\theta_0=\pi/2$ and $\theta_l=F(\theta_{l-1})$, where $\cos F(\theta)=(\sin\theta+(\pi-\theta)\cos\theta)/\pi$ (the arc-cosine kernel).
$\Fa$ is the Fatou coordinate of stage 11 ($\Fa\circ F=\Fa+1$, $\Fa(\theta)=3\pi/\theta+\frac32\log\theta+O(\theta)$, $\Fa(\pi/2)=7.7548$).
The standard normal density and distribution function are $\varphi$ and $\Phi$. Single-neuron moments of $\relu(r+\xi)$, $\xi\sim N(0,1)$:
\[
m_1(r)=\varphi(r)+r\Phi(r),\qquad m_2(r)=(1+r^2)\Phi(r)+r\varphi(r),\qquad s_1(r)^2=m_2(r)-m_1(r)^2 .
\]

\subsection{The abstract two-tier system}
\begin{definition}[two-tier system]\label{def:twotier}
A \emph{two-tier system} consists of the following.
\begin{itemize}[leftmargin=*]
\item An \emph{upper tier}: a manifold $\mathcal U$ of quasi-free (Gaussian) states.
\item A \emph{lower tier}: a finite lattice gas $\{0,1\}^n$ fibred over $\mathcal U$.
\item A \emph{down map} $\mathcal D:\mathcal U\to\mathcal P(\{0,1\}^n)$, sending a state to the conditional law of the gas.
\item An \emph{up map} $\mathcal R$, sending the joint law of gas and field to the next upper-tier state.
\end{itemize}
The \emph{closure} is the composite $\mathcal R\circ\mathcal D$ in which the gas is integrated against the quasi-free state. The
\emph{memory} is the difference between the true next state and the closure.
\end{definition}
In stages 9 and 10 the lower tier was the network and the upper tier was Eldan's localization path of the input. The tiers were
coupled from outside, through conditioning on $\cF_t$. The point of this note is that one layer of the network is already a
two-tier system of this kind.

\subsection{The cycle of one layer}
\begin{itemize}[leftmargin=*]
\item \textbf{Upper tier at layer $l$:} the quasi-free projection $(\mu_l,C_l)$ of the law of $z_l$. By positive homogeneity only rays
matter, and each neuron contributes a point $(\mu_{l,a},s_a)$ of the upper half plane. The upper tier is a cloud of $n$ points in a
hyperbolic plane (Section~\ref{sec:hyper}).
\item \textbf{Lower tier at layer $l$:} the gates $g_l\in\{0,1\}^n$. Under the quasi-free law their joint law is the dichotomized Gaussian
$g_a=\1\{r_a+\xi_a>0\}$, with $\xi\sim N(0,c)$ and $c_{ab}=C_{l,ab}/(s_as_b)$ (Section~\ref{sec:gates}).
\item \textbf{Up map:} $h_l=D_lz_l$, the field gated by the gas. Then $(m_l,\Sigma_l)$ is pushed forward,
$(\mu_{l+1},C_{l+1})=(W_{l+1}m_l,W_{l+1}\Sigma_lW_{l+1}^{\top})$.
\item \textbf{Closure:} compute $(m_l,\Sigma_l)$ from $(\mu_l,C_l)$ as if $z_l$ were Gaussian. Then
$m_{l,a}=s_am_1(r_a)$, $\Sigma_{l,aa}=s_a^2s_1(r_a)^2$, and $\Sigma_{l,ab}$ is the bivariate rectified covariance. This is the
second-order moment chain of stages 9 and 10.
\item \textbf{Memory:} the non-quasi-free part of the law of $z_l$. Its size is measured by the negentropy
$D(\mathrm{law}(z_l)\,\|\,N(\mu_l,C_l))$, and its effect on means is given by the cumulants of order $\ge3$
(Section~\ref{sec:coupling}).
\end{itemize}

\subsection{The exact split}
\begin{proposition}[Euler split; proved]\label{prop:split}
Let $f=c^{\top}h_L$, or any output that is positively homogeneous of degree one in $z_l$; every deeper output is. Put
$J_l=\partial f/\partial z_l$, a piecewise constant row vector. Then $f=\langle J_l,z_l\rangle$ pointwise, and
\[
\E f=\underbrace{\langle\E J_l,\mu_l\rangle}_{\text{upper tier}}+\underbrace{\textstyle\sum_a\Cov(J_{l,a},z_{l,a})}_{\text{lower tier}}.
\]
If $z_l$ is Gaussian, Stein's identity turns the lower-tier term into $\langle C_l,\E\nabla^2_{z_l}f\rangle$. The Hessian of the piecewise
linear $f$ is a measure on the walls of layers $\ge l$, so the lower-tier term is then a border functional. At the last layer the
split is, neuron by neuron,
\[
\E h_{L,a}=p_a\mu_a+\Cov(g_a,z_a),\qquad\text{and under the closure}\quad s_am_1(r_a)=\underbrace{s_ar_a\Phi(r_a)}_{\text{upper}}+\underbrace{s_a\varphi(r_a)}_{\text{lower}}.
\]
At $l=0$, where $z_0:=x$ and $\mu_0=0$, the upper term vanishes: $\E f=\E\langle x,\nabla f\rangle=\E\Delta f$. The whole mean is then a
lower-tier border functional, which is stage 11's border formula.
\end{proposition}
\begin{proof}
Euler's identity for degree-one homogeneous piecewise linear maps gives $f=\langle J_l,z_l\rangle$. Take expectations and write
$\E[J_az_a]=\E J_a\,\mu_a+\Cov(J_a,z_a)$. For Gaussian $z$, Stein gives $\Cov(J_a,z_a)=\sum_bC_{ab}\E\partial_bJ_a$, and
$\partial_bJ_a=\partial_a\partial_bf$. The last-layer form uses $J_{L,a}=c_ag_a$ and $\Cov(\1\{z>0\},z)=s^2p_z(0)=s\varphi(r)$ for
$z\sim N(rs,s^2)$.
\end{proof}
The same split holds with the post-activations and the Putnam--Trevi\~no pairing of stage 11. Since $f=\langle h_l,\delta_l\rangle$ with
$\delta_l=\partial f/\partial h_l$, we have $\E f=\langle m_l,\E\delta_l\rangle+\sum_i\Cov(h_{l,i},\delta_{l,i})$. The upper-tier term pairs
the two harmonic states' means, and the lower-tier term their fluctuations. Section~\ref{sec:modular} shows that the forward state
localizes with depth, while the backward state delocalizes towards the input.

\section{The internal Eldan clock}
\label{sec:clock}
\subsection{Eldan's localization seen through one row}
Eldan's localization of $\gamma$ tilts it by $\exp(\langle y,x\rangle-t|x|^2/2)$. The result is the ball $\gamma_{t,y}=N(y/(1+t),I/(1+t))$,
and along the localization path ($y_t=tX+B_t$) the centre $y_t/(1+t)$ has law $N(0,\frac t{1+t}I)$. Project onto a unit row $w$:
the ball becomes $N(\nu,\frac1{1+t})$, with standardized centre $\nu\sqrt{1+t}\sim N(0,t)$. Two independent samples from the same
localized ball have overlap $t/(1+t)+o(1)$ in high dimension. So a first layer that sees the input at localization time $t$ sees,
through every row, a unit ball at a field distributed as $N(0,t)$.

\begin{lemma}[field law; proved]\label{lem:field}
For $t\ge0$, $\E_{r\sim N(0,t)}\,\Phi(r)\Phi(-r)=\dfrac{\arccos\frac t{1+t}}{2\pi}$.
\end{lemma}
\begin{proof}
$\Phi(r)\Phi(-r)=\Pr(\xi_1<r,\ \xi_2<-r)$ for independent standard $\xi_i$. Averaging over $r$, the pair $(\xi_1-r,\xi_2+r)$ is centred
Gaussian with variances $1+t$ and correlation $-t/(1+t)$. The orthant formula gives $\frac14-\frac1{2\pi}\arcsin\frac t{1+t}=\frac1{2\pi}\arccos\frac t{1+t}$.
\end{proof}

\subsection{The internal clock}
\begin{definition}
For a fixed network, the \emph{internal clock} and \emph{internal overlap} of layer $l$ are
\[
t_l=\frac{|m_l|^2}{\Tr\Sigma_l},\qquad q_l=\frac{|m_l|^2}{\E|h_l|^2}=\frac{t_l}{1+t_l}.
\]
Since $|m_l|^2=\E\langle h_l(x),h_l(x')\rangle$ for independent inputs $x,x'$, $q_l$ is the \emph{replica overlap} of two independent inputs at depth $l$.
\end{definition}

\begin{theorem}[self-localization; proved at infinite width]\label{thm:selfloc}
Fix $l$ and let $n\to\infty$.
\begin{enumerate}[leftmargin=*]
\item $q_l\to\cos\theta_l$ and $t_l\to t_l^\infty:=\cos\theta_l/(1-\cos\theta_l)$. Equivalently $\arccos\frac{t_l}{1+t_l}=\theta_l$: the internal clock is the Eldan
time whose input overlap equals the folding angle.
\item The empirical law over $a\in[n]$ of the fields $r_{l+1,a}$ of the next layer converges to $N(0,t_l^\infty)$. For almost every row, the
law of $(z_{l+1,a}-\mu_{l+1,a})/s_a$ over $x$ converges to $N(0,1)$.
\item The mean gate variance of layer $l+1$ converges to $\frac1n\sum_ap_a q_a\to\E_{N(0,t_l)}\Phi(r)\Phi(-r)=\theta_l/2\pi$.
\item The clock obeys the closed recursion
\[
t_{l+1}^\infty=\frac{\E\,m_1(r)^2}{\E\,s_1(r)^2},\qquad r\sim N(0,t_l^\infty),
\]
and in the chart $t\mapsto\arccos\frac t{1+t}$ this recursion is the folding map $F$.
\end{enumerate}
So layer $l+1$ looks at the law of $h_l$, through each of its rows, exactly as a first layer looks at Eldan's localized input at time
$t_l$. The rows of $W_{l+1}$ play the role of localization paths, and self-averaging over neurons replaces averaging over paths.
\end{theorem}
\begin{proof}
(1) For independent $x,x'$ the pair $(h_l(x),h_l(x'))$ obeys the arc-cosine (NNGP) recursion. The input angle concentrates at $\pi/2$, so
the normalized overlap is $\cos\theta_l$ in the limit. Also $|m_l|^2=\E\langle h_l(x),h_l(x')\rangle$ because $x,x'$ are independent given
$W$.
(2) $W_{l+1}$ is independent of $(m_l,\Sigma_l)$, so $\mu_{l+1,a}=\langle w_a,m_l\rangle\sim N(0,2|m_l|^2/n)$ exactly. The quadratic form
$s_a^2=w_a^{\top}\Sigma_lw_a$ concentrates at $\frac2n\Tr\Sigma_l$ because $\|\Sigma_l\|_{op}=o(\Tr\Sigma_l)$ at fixed $l$: the largest
(collective) eigenvalue carries a share $O(1/n)$, see Proposition~\ref{prop:collective}. The ratio therefore converges in law to
$N(0,|m_l|^2/\Tr\Sigma_l)$. Gaussianity of the projected law is Sudakov's (Diaconis--Freedman's) theorem for random projections of a
thin-shell vector, and thin shell holds by the Hanin--Nica log-normal law of $|h_l|^2$ (Proposition~\ref{prop:diffusion}).
(3) follows from (2) with $p_a=\Phi(r_a)$ and Lemma~\ref{lem:field}.
(4) Under (2), $m_{l+1,a}=s\,m_1(r_a)$ and $\Sigma_{l+1,aa}=s^2s_1(r_a)^2$, which gives the ratio. For the chart, mixing over $r$ turns
$z=s(r+\xi)$ into $N(0,s^2(1+t))$, so $s^2\E m_2=s^2(1+t)/2$. Likewise $s^2\E m_1(r)^2=\E[\relu(z)\relu(z')]$ for two draws sharing
the centre, with correlation $t/(1+t)=\cos\theta$. By Cho--Saul this is $\frac{s^2(1+t)}{2\pi}(\sin\theta+(\pi-\theta)\cos\theta)$. Hence
$\frac{t'}{1+t'}=\frac{\E m_1^2}{\E m_2}=\cos F(\theta)$.
\end{proof}

\begin{remark}[what is and is not proved]
Items (1), (3) and (4) are the standard infinite-width (NNGP) statements, rephrased. The new content is the identification in (2):
the rows of the next layer sample Eldan centres at time $t_l$. The projection CLT used in (2) is rigorous given thin shell. We have not
written out rates, and at fixed $n$ the collective mode makes $s_a$ depend on $\mu_a$ (Proposition~\ref{prop:collective}).
\end{remark}

\subsection{Depth is localization time}
\begin{corollary}[one clock; proved at infinite width]\label{cor:oneclock}
Put $\Theta(t)=\arccos\frac t{1+t}$ and $\cT(t)=\Fa(\Theta(t))$. Then $\cT(t_l^\infty)=\Fa(\pi/2)+l$ exactly. Externally localizing the input to time
$t$ and then running $l$ layers gives internal time $t'$ with $\cT(t')=\cT(t)+l$. This is the stage-11 collapse theorem seen from inside:
the gates of layer $l+1$ at external time $t$ sit at the stage-11 Fatou time $\tau(l+1,t)=\cT(t)+l$.
\end{corollary}
Write $\tau_l:=\Fa(\theta_l)=7.755+l$ for the Fatou time of the representation $h_l$; it is also the Fatou time of the gates of layer $l+1$.
Szekeres' regular iteration $F^s=\Fa^{-1}(\Fa+s)$ is the unique $C^1$ flow embedding of the parabolic germ that is regular at $0$. It
defines \emph{fractional depth}, and Corollary~\ref{cor:oneclock} says that Eldan localization to time $t$ \emph{is} fractional depth
\[
s(t)=\Fa(\Theta(t))-\Fa(\pi/2)\ \sim\ 3\pi\sqrt{(1+t)/2}\qquad(t\to\infty).
\]
Exactly, $\sqrt{1+t_l}=1/(\sqrt2\sin(\theta_l/2))$. One layer is worth $\Delta\sqrt{1+t}\to\sqrt2/3\pi=0.150$ (stage 11), and
$t_l^\infty=0.47,0.98,2.13,5.04,8.72,13.18$ at $l=1,2,4,8,12,16$. The infinite-width network is therefore the time-one map of a flow
on the clock line, and the external upper tier of stages 9--10 is the same flow started earlier.

\subsection{The network cools itself}
\begin{proposition}[cooling; proved]\label{prop:cooling}
$\gamma_{t,y}\propto\exp(-(1+t)H+\langle y,x\rangle)$ with $H(x)=|x|^2/2$. So the localized ball is the Gibbs (KMS) state of $H$ at inverse
temperature $\beta=1+t$ with external field $y$, and $\hbar=1/\beta$ is a temperature. At infinite width layer $l$ is at internal
temperature
\[
T_l=\frac1{1+t_l}=1-\cos\theta_l=2\sin^2\tfrac{\theta_l}2\approx\frac{\theta_l^2}2,
\]
and its unresolved gate fraction $\theta_l/\pi=\frac2\pi\arcsin\sqrt{T_l/2}$ freezes like $\sqrt{T_l}$. The schedule $T_l\approx\frac{9\pi^2}{2\tau_l^2}$ is an
annealing schedule in Fatou time, and it is critical: the accumulated unresolved fraction $\sum_l\theta_l/\pi\approx3\log\tau$ diverges
logarithmically.
\end{proposition}

\begin{proposition}[angular and radial temperature; heuristic, consistent with data]\label{prop:Tsplit}
Write $h_l=\rho\,u$ with $\rho^2=|h_l|^2/\E|h_l|^2$, and let $V_l=\Var_x\log|h_l|^2$. Suppose $\rho$ and $u$ are independent and $\log\rho^2$
is Gaussian; independence is exact for the input's own radius, since $h_l(x)=|x|\,h_l(x/|x|)$. Then
\[
T_l=1-e^{-V_l/4}\cos\angle_l=\underbrace{(1-\cos\angle_l)}_{T^{\rm ang}_l}+\underbrace{V_l/4}_{T^{\rm rad}_l}+O(V_lT_l^{\rm ang}),
\]
where $\cos\angle_l$ is the overlap of the normalized representations, which tends to $\cos\theta_l$. The radial temperature is a
finite-width effect. It is the log-normal norm, and it does not vanish with depth: by Proposition~\ref{prop:diffusion}, $V_l$ saturates at
$44.6/n$. The two temperatures cross at $\theta^2/2=V_\infty/4$, i.e.\ at Fatou time $\tau_*\approx3\pi\sqrt{2n/44.6}=2.0\sqrt n$. Beyond
it the clock saturates at $t\approx4/V$, while the angle keeps shrinking (compare the stage-11 conjecture on the end of the parabolic
regime).
\end{proposition}
\begin{proof}[Derivation]
$q_l=|\E[\rho u]|^2/\E|h_l|^2=(\E\rho)^2\cos\angle_l$ under independence. With $\log\rho^2\sim N(-V/2,V)$, $(\E\rho)^2=e^{-V/4}$.
\end{proof}
At $n=1024$ and $l=16$ the measured temperature is $T=1/(1+11.84)=0.0779$. Proposition~\ref{prop:Tsplit} with the measured $V_{16}=0.0321$
predicts $1-e^{-0.0080}\cos\theta_{16}=0.0779$, so the radial part accounts for the whole gap between the measured clock $11.84$ and the
infinite-width $13.18$. At $l=4,8,12$ the same comparison is off by $+0.024$, $+0.010$, $-0.003$. That scatter has the size of the quenched
fluctuations of the clock. A width scan over $l=2,\dots,8$ ($n=256,1024,4096$; $6,3,2$ networks) gave mean deviations
$\cos\theta_l-q_l=+0.012,-0.007,-0.001$, which change sign and shrink with $n$: they are quenched $O(n^{-1/2})$ fluctuations of a random
clock, not a bias. The radial part is a systematic $O(1/n)$ effect. At these widths it is visible only in the deepest layers, where it
has grown to $V/4\approx0.008$.

\subsection{Measured self-localization}
A single He network ($n=1024$, $L=16$, $2\times10^5$ standard inputs, seed 5) gives the following; full output in the Verification section.
\begin{center}\small
\begin{tabular}{rcccc}\toprule
$l$ & $q_l$ meas./$\cos\theta_l$ & $t_l$ meas./$t_l^\infty$ & gate var.\ meas./$\theta_{l-1}/2\pi$ & $V_l$ meas./pred.\\\midrule
1 & 0.3183 / 0.3183 & 0.47 / 0.47 & 0.2500 / 0.2500 & 0.0069 / 0.0068\\
2 & 0.4857 / 0.4937 & 0.94 / 0.98 & 0.1994 / 0.1984 & 0.0107 / 0.0107\\
4 & 0.6545 / 0.6810 & 1.89 / 2.13 & 0.1517 / 0.1466 & 0.0172 / 0.0164\\
8 & 0.8196 / 0.8344 & 4.54 / 5.04 & 0.1009 / 0.1000 & 0.0249 / 0.0233\\
12 & 0.8933 / 0.8971 & 8.38 / 8.72 & 0.0765 / 0.0769 & 0.0292 / 0.0274\\
16 & 0.9221 / 0.9295 & 11.84 / 13.18 & 0.0614 / 0.0629 & 0.0321 / 0.0300\\
\bottomrule\end{tabular}
\end{center}
The field law of the next layer is checked in Section~\ref{sec:hyper}. The fraction of unresolved neurons ($|r|<1$) at $l=4,8,12,16$
is listed in Table~\ref{tab:fields}.

\section{The upper tier is hyperbolic}
\label{sec:hyper}
\subsection{The Fisher--Rao plane and the kink geodesic}
The Fisher--Rao metric on the univariate normal family $N(\mu,\sigma^2)$ is $ds^2=(d\mu^2+2\,d\sigma^2)/\sigma^2$. In the coordinates
$(\mu/\sqrt2,\sigma)$ it is twice the Poincar\'e metric, so the family is a hyperbolic plane of curvature $-\frac12$ (Skovgaard). We call
the vertical line $K=\{\mu=0\}$, the centred balls, the \emph{kink geodesic}: these are the laws that sit exactly on a \textsc{ReLU}
wall.
\begin{proposition}[proved]\label{prop:hyper}
\begin{enumerate}[leftmargin=*]
\item The dilations $(\mu,\sigma)\mapsto(\lambda\mu,\lambda\sigma)$ are the hyperbolic translations along $K$. Their orbits are the
equidistant curves $\{r=\mu/\sigma=\mathrm{const}\}$.
\item The distance of $N(\mu,\sigma^2)$ to $K$ is $d=\sqrt2\operatorname{asinh}(|r|/\sqrt2)$. So the field is a signed hyperbolic distance,
$r=\pm\sqrt2\sinh(d/\sqrt2)$.
\item The single-neuron rectification $\varrho:(\mu,\sigma)\mapsto(\E\relu,\operatorname{sd}\relu)=\sigma\,(m_1(r),s_1(r))$ commutes with the
translations along $K$. It is a skew product over the leaf space of equidistants: the base map is $R(r)=m_1(r)/s_1(r)$, and the
translation cocycle is $\log s_1(r)$, i.e.\ the log-scale moves along $K$ by $\log s_1(r)$.
\item $R(0)=0.6833$, $R(r)=r+o(1)$ as $r\to+\infty$, and $R(r)\sim\sqrt{\varphi(r)/2|r|}\to0^+$ as $r\to-\infty$. Hence $R(\R)\subset(0,\infty)$:
after rectification every neuron lies on the positive side of $K$.
\end{enumerate}
\end{proposition}
\begin{proof}
(1) and (2): translate to the Poincar\'e half-plane, where the distance from $(u,\sigma)$ to the imaginary axis is $\operatorname{asinh}(|u|/\sigma)$.
(3): $\relu$ is positively homogeneous. (4): use the Mills-ratio asymptotics $m_1\sim\varphi/r^2$ and $m_2\sim2\varphi/|r|^3$ as $r\to-\infty$.
\end{proof}
Item (4) is the geometric reason the clock runs forward: rectification moves the whole line of fields to one side of the kink
geodesic. The next layer's mixing then spreads the cloud again, but with a positive bias in the mean direction. The aggregate effect
is the ratio of averages $t'=\E m_1^2/\E s_1^2$, not the average of $R^2$.

\subsection{The upper tier as a cloud in the hyperbolic plane}
At layer $l$ the upper tier is the cloud $\{(\mu_{l,a},s_a)\}_{a\le n}$. By Theorem~\ref{thm:selfloc} its signed distances to $K$ satisfy
$\sqrt2\sinh(d_a/\sqrt2)\sim N(0,t_{l-1})$, so a typical neuron is at distance $d\approx\sqrt2\log\sqrt{2t}\approx\sqrt2\log l+O(1)$. The cloud
escapes from the kink geodesic logarithmically in depth. Only its second moment in the $r$-chart, the clock, enters the infinite-width
dynamics, and the dilation coordinate along $K$ drops out by homogeneity.
\begin{proposition}[criticality and diffusion along $K$; drift proved, diffusion as in Hanin--Nica]\label{prop:diffusion}
The He variance $2/n$ is the unique scale for which $\E|h_l|^2$ is constant in $l$: the barycentric log-scale has zero drift along
$K$. At finite width, the input-dependent log-norm performs a random walk along $K$. For a fixed network its across-input increments
have variance
\[
\Var_x\Big(\log\frac{|h_l|^2}{|h_{l-1}|^2}\Big)=\frac4n\Big(\frac32-\frac{J_2(\theta_{l-1})}{2\pi}\Big),\qquad J_2(\theta)=3\sin\theta\cos\theta+(\pi-\theta)(1+2\cos^2\theta),
\]
with $V_0=2/n$. Since $\frac32-\frac{J_2(\theta)}{2\pi}=\theta^2+O(\theta^4)$, the increments are summable in the parabolic regime and
\[
nV_{16}=30.76,\qquad nV_\infty=44.60,
\]
against $2+5l=82$ at $l=16$ for the joint (weights and input) variance of Hanin--Nica. The input-dependent part of the norm is made in
the first layers: deep layers act on all inputs alike, because two inputs at angle $\theta$ disagree on only a fraction $\theta/\pi$ of
the gates.
\end{proposition}
\begin{proof}[Derivation]
Per layer the norm ratio is
\[
|h_l|^2/|h_{l-1}|^2\approx\tfrac2n\textstyle\sum_a\relu(\zeta_a)^2,\qquad\zeta_a\ \text{standard},
\]
with variance $\frac4n(\E\relu^4-\frac14)=5/n$.
The across-input variance is this minus the covariance between two independent inputs, whose $\zeta$'s are at angle $\theta_{l-1}$. That
covariance is
\[
\tfrac4n\big(\E[\relu(\zeta)^2\relu(\zeta')^2]-\tfrac14\big)=\tfrac4n\big(\tfrac{J_2(\theta)}{2\pi}-\tfrac14\big)
\]
by the degree-two arc-cosine kernel.
Expanding $J_2$ at $0$ gives $3\pi-2\pi\theta^2+O(\theta^4)$.
\end{proof}

\begin{proposition}[the collective mode; heuristic, measured]\label{prop:collective}
The covariance $\Sigma_l$ has a collective direction along $m_l$, carrying the share $\approx t_lV_l/4$ of the trace. This is the norm
fluctuation projected on the mean: $\langle\hat m,h\rangle\approx|m|(1+\delta)$ with $\Var\delta\approx V/4$, i.e.\ the diffusion along $K$.
Through the next layer it makes the scale depend on the mean, $s_a^2\approx s_{\rm bulk}^2+\epsilon\mu_a^2$ with $\epsilon\approx V/4$. The
field law is therefore compressed,
\[
r_a=\frac{\varrho_a}{\sqrt{1+\epsilon\varrho_a^2}},\qquad \varrho_a\sim N(0,t_{\rm bulk}),\qquad |r_a|<\epsilon^{-1/2}\approx2/\sqrt V .
\]
In the hyperbolic picture the cloud cannot go further than $\sqrt2\operatorname{asinh}\sqrt{2/V}$ from the kink geodesic. The collective
share becomes $O(1)$ when $t_lV_l/4\sim1$, which is again $\tau_*\approx2\sqrt n$ (Proposition~\ref{prop:Tsplit}).
\end{proposition}
\begin{table}[h]\centering\small
\begin{tabular}{rccccccc}\toprule
$l$ & $\Var r$ & $t_{l-1}$ meas. & compressed & $t_{l-1}^\infty$ & $\Pr(|r|<1)$ & coll./$tV/4$ & Frob.\\\midrule
4 & 1.38 & 1.43 & 1.41 & 1.53 & 0.615 / 0.581 & 0.008 / 0.005 & 0.034\\
8 & 3.79 & 3.74 & 3.57 & 4.24 & 0.362 / 0.373 & 0.026 / 0.022 & 0.122\\
12 & 7.07 & 7.51 & 6.83 & 7.73 & 0.286 / 0.281 & 0.058 / 0.053 & 0.316\\
16 & 10.14 & 10.93 & 9.52 & 11.99 & 0.226 / 0.227 & 0.093 / 0.086 & 0.469\\
\bottomrule\end{tabular}
\caption{Field law of layer $l$ ($n=1024$, one network). $\Pr(|r|<1)$ is measured against $\operatorname{erf}(1/\sqrt{2t^\infty_{l-1}})$, and
``coll.'' is the collective trace share against $tV/4$. ``Compressed'' is $\E[\varrho^2/(1+\epsilon\varrho^2)]$ with $t_{\rm bulk}$ and $\epsilon$ taken from the
measured clock and collective share. The last column (``Frob.'') is stage 9's metric, the Frobenius share of the top eigenvalue of the
off-diagonal covariance.}\label{tab:fields}
\end{table}
Table~\ref{tab:fields} supports the picture. The unresolved fraction $\Pr(|r|<1)$ matches the infinite-width $\operatorname{erf}(1/\sqrt{2t})$ to
$0.001$ at the deepest layers, because compression acts on large $|r|$. The measured collective trace share tracks $tV/4$
($0.093$ against $0.086$ at $l=16$). At $l=12$ and $16$ the field variance lies below the measured clock, as compression predicts, but the one-parameter model
over-predicts the effect: at $l=16$ it gives $9.52$, against $10.14$ measured and $10.93$ uncompressed. At $l=8$ the variance exceeds the
measured clock, which is within the quenched scatter. The over-prediction disappears once only a fraction $\eta\approx0.71$ of $V$ is
counted as a common scale: Proposition~\ref{prop:skewlaw} measures this efficiency independently from the skewness, and it then gives
$7.11$ and $10.09$ against $7.07$ and $10.14$. The stage-9 statement that ``37--48\% of the off-diagonal covariance
energy is collective'' is the Frobenius share in the last column. It reaches $0.47$ at $l=16$ although the collective \emph{trace} share is
only $0.09$: a rank-one mode of modest trace dominates the off-diagonal Frobenius norm, because the bulk off-diagonal entries are
random and incoherent.

\section{The lower tier: the gate gas and the wall-density law}
\label{sec:gates}
\subsection{The gate gas}
For a fixed network the \emph{gate gas} of layer $l$ is the law of $g_l(x)\in\{0,1\}^n$, $x\sim\gamma$. Under the quasi-free law of $z_l$ it is
the dichotomized Gaussian of Emrich--Piedmonte and Macke et al.: $g_a=\1\{r_a+\xi_a>0\}$ with $\xi\sim N(0,c)$. The one-body marginals are
$p_a=\Phi(r_a)$. The two-body marginals are orthant probabilities $\Phi_2(r_a,r_b;c_{ab})$.

\begin{proposition}[Ising couplings; proved to first order in $c$]\label{prop:ising}
The pair coupling of the gate gas (the log odds ratio of the $2\times2$ table) is
\[
J_{ab}=\log\frac{P_{11}P_{00}}{P_{10}P_{01}}=\frac{c_{ab}\,\varphi(r_a)\varphi(r_b)}{p_aq_a\,p_bq_b}+O(c_{ab}^2).
\]
\end{proposition}
\begin{proof}
$\partial_c\Phi_2(r_a,r_b;c)|_{c=0}=\varphi(r_a)\varphi(r_b)$. A covariance $\epsilon$ in a table with margins $p_a,p_b$ changes the log odds ratio by
$\epsilon(\frac1{p_ap_b}+\frac1{p_aq_b}+\frac1{q_ap_b}+\frac1{q_aq_b})=\epsilon/(p_aq_ap_bq_b)$.
\end{proof}
The correlations $c_{ab}$ have two parts.
\begin{itemize}[leftmargin=*]
\item A \emph{bulk} part of size $O(n^{-1/2})$ with random signs, which is Sherrington--Kirkpatrick-like.
\item A \emph{collective} part $c^{\rm coll}_{ab}\approx\epsilon\,r_ar_b$ coming from the norm mode (Proposition~\ref{prop:collective}).
\end{itemize}
By Proposition~\ref{prop:ising} the collective part is a Mattis coupling, $J^{\rm coll}_{ab}=\epsilon\,\xi_a\xi_b$, with stored pattern
\[
\xi_a=\frac{r_a\varphi(r_a)}{p_aq_a}\ \approx\ r_a|r_a|\quad(|r_a|\gg1).
\]
The pattern the gate gas stores is the field itself, that is, the representation's own mean. Its Mattis order parameter is not
spontaneous. It is the input's radial coordinate $\delta(x)$, a latent variable shared by all neurons. Given $\delta$, the gates are
conditionally independent up to the bulk couplings, which is a de Finetti form. Resolved neurons ($|r|\gg1$) have large couplings
per unit correlation but tiny variance $p_aq_a$; the gas is thermally active only near its ``Fermi level'' $r=0$.

\begin{remark}[two statistics; analogy]
The field $z$ is a quasi-free boson field (a Gaussian, i.e.\ a quasi-free state on a CCR algebra). The gates are fermionic occupation
numbers ($g^2=g$). In the independent-gate approximation their state is the product (gauge-invariant quasi-free) state with diagonal
one-particle density $\diag(\Phi(r_a))$, and the occupation is a probit-smeared Fermi function of the ``energy'' $-r_a$. The
\textsc{ReLU} coupling $h=g\,z$ is a boson--fermion (Yukawa-type) vertex with the occupation slaved to the field. The Euler split
(Proposition~\ref{prop:split}) is its Wick decomposition, $\E[gz]=\E g\,\E z+\Cov(g,z)$, and the closure is the corresponding
mean-field (Hartree--Fock--Bogoliubov) decoupling. Depth lowers the density of states at the Fermi level as $1/\tau$
(Theorem~\ref{thm:walldensity}).
\end{remark}

\subsection{The wall-density law}
\begin{theorem}[wall-density law; proved]\label{thm:walldensity}
Let $Q$ be integrable with $\int r^2|Q|<\infty$, and $\Theta=\arccos\frac t{1+t}$. Then
\[
\E_{r\sim N(0,t)}Q(r)=\Big(\int Q\Big)\frac{\sin(\Theta/2)}{\sqrt\pi}\Big(1-\frac{m_2(Q)-1}{2t}+O(t^{-2})\Big),\qquad m_2(Q)=\frac{\int r^2Q}{\int Q},
\]
and $\sin(\Theta/2)/\sqrt\pi=1/\sqrt{2\pi(1+t)}$ exactly. At layer $l+1$, where $\Theta=\theta_l$, every wall-localized one-body quantity therefore equals
\[
\Big(\int Q\Big)\frac{\sin(\theta_l/2)}{\sqrt\pi}\approx\Big(\int Q\Big)\frac{\theta_l}{2\sqrt\pi}\approx\frac{3\sqrt\pi}2\Big(\int Q\Big)\frac1{\tau_l}=2.659\Big(\int Q\Big)\frac1{\tau_l}.
\]
\end{theorem}
\begin{proof}
Expand $e^{-r^2/2t}/\sqrt{2\pi t}=(2\pi t)^{-1/2}(1-r^2/2t+O(t^{-2}r^4))$ and $(2\pi(1+t))^{-1/2}=(2\pi t)^{-1/2}(1-1/2t+O(t^{-2}))$. Use
$1+t=1/(2\sin^2(\Theta/2))$.
\end{proof}
\begin{center}\small
\begin{tabular}{lccl}\toprule
quantity $Q$ & $\int Q$ & law at layer $l+1$ & check\\\midrule
wall density $\varphi(r)$ & $1$ & $\sin(\theta_l/2)/\sqrt\pi$ (exact) & $t=13$: $0.10662181$ both\\
gate variance $\Phi(r)\Phi(-r)$ & $1/\sqrt\pi$ & $\theta_l/2\pi+O(\theta^3)$ & $t=100$: $0.022415$ vs $0.022396$\\
third-cumulant source $k_3(r)^2$ & $0.2700$ & $0.2700\,\sin(\theta_l/2)/\sqrt\pi$ & Section~\ref{sec:coupling}\\\bottomrule
\end{tabular}
\end{center}
\begin{corollary}[the harmonic budget; proved at infinite width]\label{cor:budget}
The activity of the lower tier decays harmonically in Fatou time, and its total over depth grows like $2.659(\int Q)\log\tau$. Each
$e$-fold of Fatou time costs the same: the lower tier, and with it the memory it gives birth to, is spread log-uniformly over depth,
not concentrated in the deepest layers. For the gate variance the total over $16$ layers is $\sum_{l\le16}\theta_{l-1}/\pi=3.658$.
\end{corollary}

\begin{proposition}[Hermite cancellation; proved]\label{prop:hermite}
$\E_{r\sim N(0,t)}[\He_k(r)\varphi(r)]=0$ for odd $k$, and for $k=2j$
\[
\E_{r\sim N(0,t)}[\He_{2j}(r)\varphi(r)]=\frac{(-1)^j(2j-1)!!}{\sqrt{2\pi}}\,(1+t)^{-j-1/2}.
\]
\end{proposition}
\begin{proof}
$\varphi(r)\,\varphi_t(r)=(2\pi(1+t))^{-1/2}\varphi_{\sigma^2}(r)$ with $\sigma^2=t/(1+t)$. For $Y\sim N(0,\sigma^2)$, $\E\He_{2j}(Y)=(\sigma^2-1)^j(2j-1)!!$, from
the generating function $\E e^{sY-s^2/2}=e^{s^2(\sigma^2-1)/2}$.
\end{proof}
The marginal memory of a neuron, i.e.\ the error of the one-neuron closure, is by the Gram--Charlier form of the Edgeworth series
(Section~\ref{sec:coupling})
\[
\frac{\E\relu(z_a)-s_am_1(r_a)}{s_a}=\Big(-\frac{\gamma_{1,a}}6\He_1(r_a)+\frac{\gamma_{2,a}}{24}\He_2(r_a)+\frac{\gamma_{1,a}^2}{72}\He_4(r_a)\Big)\varphi(r_a)+\dots,
\]
where $\gamma_1,\gamma_2$ are the standardized skewness and excess kurtosis of $z_a$. Every term is wall-localized. By
Proposition~\ref{prop:hermite}, its \emph{average} over neurons is suppressed by extra powers of $T=1/(1+t)$ when the $\gamma$'s are
uncorrelated with $r$. Its \emph{mean square}, the quantity an MSE score sees, obeys the wall-density law with $\int(\He_k\varphi)^2$. The
marginal memory is therefore a sign-alternating wall-localized field over neurons. It nearly cancels in neuron averages, but not in
the MSE.

On the width-1024 network the measured marginal skewness and kurtosis of the pre-activations grow with depth: mean $|\gamma_1|=0.043$,
$0.070$, $0.095$, and mean $|\gamma_2|=0.041$, $0.055$, $0.066$, at $l=8,12,16$. The sampling noise is about $0.004$ and $0.009$. Their
source is the collective mode, not the birth law (Proposition~\ref{prop:skewlaw}). With the empirical
cumulants, the three Gram--Charlier wall terms reproduce the empirical one-neuron closure error neuron by neuron: the correlation is
$0.983$, $0.993$ and $0.996$ at $l=8,12,16$, and the residual is $\le1.7\times10^{-4}$ in every $|r|$-bin. The error vanishes for $|r|\ge3$ (rms
$4\times10^{-5}$) and peaks at $1\le|r|<2$ (rms $1.9\times10^{-3}$ at $l=16$). Both sides are computed from the same $2\times10^5$ samples, so
this checks the expansion and the wall-localized profile. It does not fix the absolute size of the true memory, which is comparable to
the Monte Carlo floor ($1.4$--$1.6\times10^{-3}$).

\section{Coupling the tiers: the memory as a stratified cumulant series}
\label{sec:coupling}
\subsection{The Edgeworth series of a suffix}
Fix a layer $l$ and let $S:\R^n\to\R$ be the suffix from $z_l$ to the output; it is piecewise linear and positively homogeneous. Let
$\kappa_k$ be the cumulant tensors of $z_l$ and $\E_G$ the expectation under the quasi-free projection $N(\mu_l,C_l)$. The density of $z_l$
is formally $\exp\big(\sum_{k\ge3}(-1)^k\kappa_k\cdot\partial^k/k!\big)$ applied to the Gaussian density. Integrating by parts gives the
formal Edgeworth series (McCullagh)
\[
\E f=\E_GS+\tfrac16\langle\kappa_3,\E_G\nabla^3S\rangle+\tfrac1{24}\langle\kappa_4,\E_G\nabla^4S\rangle+\tfrac1{72}\langle\kappa_3^{\otimes2},\E_G\nabla^6S\rangle+\cdots
\]
The first term is the closure (upper tier composed with lower tier). The rest is the memory. Equivalently, by the cumulant expansion
$\E[z_aG(z)]=\sum_k\frac1{k!}\kappa_{k+1}(a,\cdot)\,\E\nabla^kG$ (Lytova--Pastur), the lower-tier term of the Euler split is
$\sum_a\Cov(J_a,z_a)=\sum_{k\ge1}\frac1{k!}\langle\kappa_{k+1},\E\nabla^{k+1}f\rangle$. Its $k=1$ term is the Stein (Gaussian) part of
Proposition~\ref{prop:split}.

For piecewise linear $S$ each $\nabla^kS$, $k\ge2$, is a distribution. Every derivative that hits a gate produces a $\delta$ (or a
derivative of $\delta$) on that neuron's wall. So each term of $\nabla^kS$ is supported on an intersection of walls, a stratum of the
wall arrangement, whose codimension is the number of distinct $\delta$ factors.

\subsection{Which strata carry memory}
\begin{proposition}[chain support; proved]\label{prop:chain}
Every stratum that carries a term of $\nabla^kS$ is a \emph{chain}: a set of walls in which every wall except the deepest lies in the
input cone of a deeper wall of the set. Strata spanned by walls of a single layer (normal crossings) carry no term. In particular, if
the suffix is one layer with a linear readout, $S=\sum_bW_b\relu(z_b)$, then all mixed derivatives vanish and the memory is purely
marginal: $\E f-\E_GS=\sum_bW_b(\E\relu(z_b)-\E_G\relu(z_b))$.
\end{proposition}
\begin{proof}
A gate of neuron $a$ depends on $z_b$ only through the input cone of $a$. Differentiating $\sigma'(y_a)\prod(\cdots)$ in $z_b$ and $z_c$ for two
neurons of one layer produces a $\delta$ on $z_b$ or $z_c$ only together with a derivative of a deeper gate that both feed. So the two
shallow walls never appear without a deeper wall in their common output cone. Induct over layers.
\end{proof}

\begin{theorem}[stratified third order; proved]\label{thm:strata}
Consider a two-layer suffix $S(z)=\sigma(y)$ with $y=\sum_bW_b\sigma(z_b)$, with $\sigma$ smooth and $g=W\odot\sigma'(z)$ the pulled-back deep normal,
\[
\langle\kappa,\nabla^3S\rangle=\underbrace{\sigma'''(y)\,\kappa(g,g,g)}_{\text{deep wall}}
+\underbrace{3\,\sigma''(y)\sum_cW_c\,\sigma''(z_c)\,\kappa(g,e_c,e_c)}_{\text{bent crossings}}
+\underbrace{\sigma'(y)\sum_cW_c\,\sigma'''(z_c)\,\kappa_{ccc}}_{\text{hub layer}}
\]
for every symmetric $3$-tensor $\kappa$. In the \textsc{ReLU} limit ($\sigma'=\1_{>0}$, $\sigma''=\delta$, $\sigma'''=\delta'$) the three terms live on:
\begin{itemize}[leftmargin=*]
\item the deep wall $\{y=0\}$ (codimension $1$), where they pair with the third cumulant of the linearized deep field $\langle g,z\rangle$;
\item the bent crossings $\{y=0\}\cap\{z_c=0\}$ (codimension $2$), where they pair with the two-body cumulant
$\kappa(\langle g,z\rangle,z_c,z_c)$ and are weighted by the bending $T_c=W_c$ of stage 11;
\item the shallow walls $\{z_c=0\}$ (codimension $1$), where they pair with the one-body cumulants $\kappa_3(z_c)$ and are weighted by the deep
gate.
\end{itemize}
\end{theorem}
\begin{proof}
Chain rule: $\partial_aS=\sigma'(y)g_a$; $\partial_a\partial_bS=\sigma''(y)g_ag_b+\delta_{ab}\sigma'(y)W_a\sigma''(z_a)$; one more derivative gives the three
groups, and the middle one appears in three index positions. Contract with symmetric $\kappa$.
\end{proof}
\begin{center}\small
\begin{tabular}{lcccccc}\toprule
softplus $\beta$ & Hermite IBP & deep wall & bent crossings & hub layer & sum & difference (s.e.)\\\midrule
$5$ & $0.1242$ & $-0.0159$ & $-0.1514$ & $0.2896$ & $0.1224$ & $0.0018\ (0.0028)$\\
$20$ & $0.1677$ & $0.2421$ & $-0.5730$ & $0.5007$ & $0.1698$ & $-0.0022\ (0.0039)$\\\bottomrule
\end{tabular}
\end{center}
The check uses a $3$-neuron Gaussian $z$ with a random symmetric $\kappa$ and $2\times10^7$ samples. The left side is
$\E[S(z)\He_\kappa(z)]$ by Hermite integration by parts; the right side is the three strata. They agree within one standard error.
The strata cancel substantially, so no single one dominates. A truncation that keeps one stratum and drops the others is not
justified by this structure.

\begin{corollary}[Kikuchi level $=$ codimension; proved at third order, by Proposition~\ref{prop:chain} in general]\label{cor:kikuchi}
The cumulant that enters a stratum is the joint cumulant of exactly the fields whose walls meet there: deep and shallow neurons along a
chain, with the deep field linearized. One-body cumulants live on codimension-one walls, two-body cumulants on codimension-two bent
crossings, and $j$-body cumulants on codimension-$j$ chains. The cluster (Kikuchi) level of the memory equals the codimension of its
border. The Gaussian closure is exact on normal-crossing arrangements. Two-body memory lives only where stage 11 found the non-abelian
isotropy $\aff(1)$.
\end{corollary}
This places the stage-9 decomposition: its hub-diagonal pairing is the hub layer, its coherent (mean-field) memory is the deep-wall
term, and the bent-crossing term is new.

\subsection{Birth of the third cumulant}
\begin{proposition}[birth law; proved at infinite width]\label{prop:birth}
For $z\sim N(rs,s^2)$, $\kappa_3(\relu z)=s^3k_3(r)$ with $k_3=m_3-3m_1m_2+2m_1^3$, $m_3(r)=(r^3+3r)\Phi(r)+(r^2+2)\varphi(r)$, and
$k_3(0)=c_3=(\pi+2)/(2\pi)^{3/2}=0.32646$. The mean-square source of one-body third cumulants at layer $l+1$ is
\[
B_l=\E_{N(0,t_l)}k_3(r)^2=0.2700\,\frac{\sin(\theta_l/2)}{\sqrt\pi}\big(1+O(\theta_l^2)\big)\approx\frac{0.718}{\tau_l}.
\]
\end{proposition}
Measured against $0.2700/\sqrt{2\pi t}$: $0.0802/0.1077$, $0.0493/0.0539$, $0.0290/0.0299$, $0.0151/0.0152$ at $t=1,4,13,50$. By the
harmonic budget the total one-body third-cumulant source over depth grows like $0.718\log\tau$, at the same rate for every $e$-fold of
Fatou time.

\begin{remark}[which route the measured skewness takes]
Transported through $W_{l+1}$ with random signs, the one-body source alone gives a standardized skewness of order $n^{-1}$. That is far
below the measured $0.095$ at $l=16$. The measured skewness must therefore come from many-body cumulants, and the natural candidate
is the collective mode. If $h=(1+\delta)\tilde h$ with $\Var\delta=v\approx V/4$ independent of the bulk, then $z_a-\mu_a=\delta\mu_a+s\xi_a(1+\delta)$,
which gives
\[
\kappa_3=6\mu_as^2v,\qquad\kappa_4=12s^4v+O(v^2),\qquad\gamma_1\approx6v\,r_a\approx\tfrac32V\,r_a,\qquad\gamma_2\approx12v\approx3V.
\]
So skewness should be linear in the field, and excess kurtosis flat in it, both proportional to $V$.
\end{remark}
\begin{proposition}[collective-mode cumulant law; measured]\label{prop:skewlaw}
On the width-1024 network the marginal cumulants of the pre-activations follow this law in form. At $l=8,12,16$ the regression of $\gamma_1$ on $r$ over the $1024$ neurons has slopes $0.0252$, $0.0302$, $0.0336$, with $R^2=0.86$, $0.90$,
$0.91$ and intercepts below $10^{-4}$. The predicted slopes $\frac32V_{l-1}$ are $0.0350$, $0.0421$, $0.0474$, a constant ratio
$\eta=0.72$, $0.72$, $0.71$. The excess kurtosis is flat in $r$ ($R^2\le0.11$ for a fit in $r^2$), at levels $0.040$, $0.053$, $0.059$ against
$3V_{l-1}=0.070$, $0.084$, $0.095$. The same efficiency $\eta$, inserted in the compression of the field law
(Proposition~\ref{prop:collective}), gives $\Var r=3.64$, $7.11$, $10.09$ against the measured $3.79$, $7.07$, $10.14$. This removes the
over-prediction in Table~\ref{tab:fields} at $l=12,16$. A fraction $\eta\approx0.7$ of the across-input log-norm variance acts as a common
random scale on the next layer's pre-activations.

\end{proposition}
The law identifies the dominant finite-width source of marginal non-Gaussianity at $L=16$, $n=1024$. It is not the one-body birth at
the walls. It is the radial diffusion along the kink geodesic, acting as a common random scale on the bulk noise. Because
\textsc{ReLU} commutes with that scale, this part of the memory is a one-dimensional mixture rather than a genuine many-body effect
(Section~\ref{sec:synthesis}, item 6).

\section{Modular structure: KMS states, centres and the two clocks}
\label{sec:modular}
\subsection{The layer state as a KMS state on the fibre relation}
Let $\mathcal R_l=\{(x,x'):h_l(x)=h_l(x')\}$ be the fibre relation of layer $l$. Its fibres are polyhedral; for $l=1$ and invertible $W_1$
the fibre through $x$ is an orthant of dimension equal to the number of closed gates. Equip it with Lebesgue measure on the fibres
(the Haar system) and the real cocycle $c(x,x')=H(x)-H(x')$, $H=|x|^2/2$. The cocycle generates the one-parameter group
$\sigma_s(k)(x,x')=e^{is\,c(x,x')}k(x,x')$ on kernels $k$ on $\mathcal R_l$.
\begin{proposition}[KMS states and the centre; proved in the measured setting]\label{prop:kms}
\begin{enumerate}[leftmargin=*]
\item The state $\omega_\mu$ of a measure $\mu$, given by $\omega_\mu(k)=\int k(x,x)\,d\mu(x)$, is $(\sigma,\beta)$-KMS iff the conditional measures of $\mu$ on the fibres are
$\propto e^{-\beta H}dx$ (Renault's correspondence, Feldman--Moore in the measured case). So the KMS$_\beta$ states are parametrized by
probability measures on the quotient $h_l(\R^n)$, and the extremal ones are the fibre Gibbs states $\omega_h$, one for each value $h$
of the representation.
\item The Gaussian state $\omega_\gamma$ is KMS$_1$, and its extremal decomposition is the law of the representation:
$\omega_\gamma=\int\omega_h\,d\mu_l(h)$ with $\mu_l=(h_l)_*\gamma$.
\item The centre of the von Neumann algebra of $\mathcal R_l$ is $\cZ_l=L^\infty(h_l)$, the functions of the representation. The
$\omega_\gamma$-preserving conditional expectation onto it is the posterior $E_l=\E[\,\cdot\,|h_l]$.
\item Eldan's localized states $\gamma_{t,y}$ are KMS at $\beta=1+t$ for the perturbed cocycle $c-\frac1\beta\langle y,x-x'\rangle$. At infinite
width layer $l$ is at $\beta_l=1+t_l$ (Proposition~\ref{prop:cooling}).
\end{enumerate}
\end{proposition}
\begin{proof}
(1) The KMS condition for $\sigma$ on the groupoid algebra says that the Radon--Nikodym cocycle of $\mu$ along $\mathcal R_l$, relative to the
Haar system, is $e^{-\beta c}$. That is the stated property of the conditionals, and extremality is ergodicity, i.e.\ a single fibre.
(2) is disintegration of $\gamma$ along $h_l$. (3) $\mathcal R_l$ has a proper (smooth) quotient, so its von Neumann algebra is the direct
integral of $B(L^2(\text{fibre}))$ over $h_l(\R^n)$, and the centre is the diagonal. (4) is Proposition~\ref{prop:cooling}.
\end{proof}
Since $h_{l+1}$ is a function of $h_l$, $\mathcal R_l\subset\mathcal R_{l+1}$. The algebras increase, and the centres form a decreasing tower
\[
L^\infty(\gamma)\supset\cZ_1\supset\cZ_2\supset\cdots\supset\cZ_L\ni f,\qquad E_{l+1}E_l=E_{l+1}.
\]
The mean is $\omega_\gamma(f)$ for $f$ in the smallest centre. This is ``localization in the commutant'' (stage 9) realized inside the
network: each layer is a central decomposition, and the representation at depth $l$ \emph{is} the spectrum of its centre. The forward
pass computes, for one input, the extremal KMS component it belongs to.

\begin{remark}[what the modular picture adds]
The upper tier of stages 9--10 (the localized balls $\gamma_{t,y}$) and the network's own tower (the centres $\cZ_l$) are two families of
states for one modular Hamiltonian $H$. The external family varies $\beta$ at fixed algebra. The internal family varies the algebra at
fixed $\beta=1$, and by Theorem~\ref{thm:selfloc} its effect on the next layer equals a change of $\beta$ to $1+t_l$. Cooling (raising
$\beta$) and coarse-graining (shrinking the centre) are thus interchangeable on the next layer, to leading order at infinite width. This
is the operator-algebraic form of ``depth is localization time''.
\end{remark}

\subsection{Two harmonic states, two clocks}
Stage 11 identified the forward activations and the backward sensitivities $\delta_l=\partial f/\partial h_l$ as the two harmonic states of
the network's Bratteli diagram, with $\langle h_l,\delta_l\rangle=f$ at every layer. Their clocks run in opposite directions.
\begin{proposition}[backward clock; infinite width, gradient independence]\label{prop:backclock}
For $f=c^{\top}h_L$, with readout $c$ independent of the weights,
\[
\frac{|\E\delta_l|^2}{\E|\delta_l|^2}\ \longrightarrow\ \prod_{k=l+1}^{L}\Big(1-\frac{\theta_{k-1}}\pi\Big)\ \approx\ \Big(\frac{\tau_l}{\tau_L}\Big)^3 .
\]
The forward state localizes with depth ($q_l=\cos\theta_l\uparrow1$), while the backward state delocalizes as it travels towards the input,
with a parabolic power law of exponent $3=1/(\pi a)$, where $F(\theta)=\theta-a\theta^2+\dots$ and $a=1/3\pi$.
\end{proposition}
\begin{proof}[Derivation]
$\E_W\langle\delta_{k-1}(x),\delta_{k-1}(x')\rangle=2\Pr(g_k=g'_k=1)\langle\delta_k,\delta'_k\rangle=(1-\theta_{k-1}/\pi)\langle\delta_k,\delta'_k\rangle$, while
$\E|\delta_{k-1}|^2=\E|\delta_k|^2$ at criticality. This uses the gradient-independence ansatz, justified by Yang's tensor programs for a
readout independent of the weights. Then $\theta_{k-1}/\pi=3/\tau+O(\tau^{-2}\log\tau)$, so the product is $\exp(-3\sum1/\tau)\approx(\tau_l/\tau_L)^3$.
\end{proof}
On the width-1024 network the measured backward clock is $0.0108$, $0.0354$, $0.0800$, $0.2489$, $0.5571$, $0.8798$ at $l=0,2,4,8,12,15$. The
infinite-width products are $0.0133$, $0.0442$, $0.0941$, $0.2663$, $0.5589$, $0.8743$. Agreement is close near the output, and the measured state
delocalizes about $20\%$ faster near the input, where finite-width errors compound. Far out, the power law is accurate: the product
over $l=1000..10^4$ is $0.00104$, against $(\Fa(\theta_{1000})/\Fa(\theta_{10^4}))^3=0.00102$.

The Euler split at layer $l$ pairs the two states: $\E f=\langle m_l,\E\delta_l\rangle+\sum_i\Cov(h_{l,i},\delta_{l,i})$. The upper-tier
(mean--mean) term needs both states localized. Near the output the forward state is localized and the backward state is trivially so.
Near the input the backward state has delocalized to $1\%$ of its energy, so the mean there is almost entirely lower tier, which is
stage 11's border formula at $l=0$. In Montanari's correspondence between stochastic localization and diffusion models, the forward
(localizing) clock and the backward (delocalizing) clock are the two time directions of one process (an analogy only).

\section{What the realisation unlocks}
\label{sec:synthesis}
\subsection{The realisation}
The two-tier framework of stages 9--11 is not imposed on the network. The network realizes it, layer by layer.
\begin{center}\small
\begin{tabular}{>{\raggedright\arraybackslash}p{4.2cm}>{\raggedright\arraybackslash}p{5.1cm}>{\raggedright\arraybackslash}p{5.3cm}}\toprule
 & external (stages 9--10) & internal (this note)\\\midrule
upper tier & Eldan balls $\gamma_{t,y}$ & quasi-free projection $(\mu_l,C_l)$; a cloud in the Fisher--Rao plane\\
clock & localization time $t$ & $t_l=|m_l|^2/\Tr\Sigma_l=\cos\theta_l/(1-\cos\theta_l)$\\
paths & Brownian localization paths & rows of $W_{l+1}$ (self-averaging)\\
temperature & $1/(1+t)$ & $1-\cos\theta_l$ (+ radial $V_l/4$ at finite width)\\
lower tier & gates conditioned on $\cF_t$ & the gate gas: dichotomized Gaussian, Mattis pattern $=$ field\\
coupling & conditioning & Euler split; Edgeworth series stratified by codimension\\
algebra & conditional expectation on $\cF_t$ & centres $\cZ_l=L^\infty(h_l)$ of the fibre relations; KMS decomposition\\
renormalization time & Fatou time $\tau(l,t)$ & $\tau_l=\Fa(\theta_l)=7.755+l$; fractional depth\\\bottomrule
\end{tabular}
\end{center}
Corollary~\ref{cor:oneclock} makes the dictionary exact at infinite width: the external upper tier is the internal one started at
fractional depth $s(t)$.

\subsection{What it unlocks}
\begin{enumerate}[leftmargin=*]
\item \textbf{One flow for all depths.} The infinite-width network is the time-one map of Szekeres' flow on the clock line, and
localization is fractional depth. A quantity calibrated at one $(l,t)$ transfers to every $(l',t')$ with the same Fatou time. Stage 11
stated this as an open problem (``chains in Fatou time''). At infinite width it is now a theorem for every quantity determined by the
pair angle. Departures from it are finite-width signals by definition.
\item \textbf{A temperature budget.} $T_l=T^{\rm ang}_l+T^{\rm rad}_l$. The angular part is the parabolic, infinite-width cooling
$\approx9\pi^2/2\tau^2$. The radial part $V_l/4$ is the log-normal norm, which saturates at $44.6/(4n)$. For the competition
networks ($n=1024$, $L=16$), $\tau_{16}=23.8$ is far below $\tau_*\approx2\sqrt n=64$, so they are angular-dominated. Yet at $l=16$ the
radial part ($0.008$) is already $11\%$ of the angular one ($0.0705$), and it accounts for the gap between the measured and the
infinite-width clock.
\item \textbf{A harmonic budget for the lower tier.} The layer average of each wall-localized one-body quantity $Q$ is $(\int Q)\sin(\theta_l/2)/\sqrt\pi$, which is about $2.66(\int Q)/\tau_l$.
The memory is born at rate $0.718/\tau$ and accumulates like $\log\tau$. A correction scheme must therefore treat every layer, with
weight $\propto1/\tau_l$. Truncating to the last few layers is not justified.
\item \textbf{The memory is a sum over bent strata.} The Edgeworth corrections decompose over chain strata of the wall arrangement,
with Kikuchi level equal to codimension. Normal crossings carry nothing, and pair memory lives only on bent crossings (the $\aff(1)$
points). This is an algorithmic skeleton: one-body corrections on single walls (hub and deep-wall terms, both already present in
stage 9), plus a pair correction on each (deep, shallow) bent crossing. Because the strata cancel, all three must be kept together.
\item \textbf{Hermite cancellation.} The neuron-averaged marginal memory is suppressed by powers of $T_l$, but its mean square is not.
Averaged diagnostics (for example the mean of the output vector) understate the memory that an MSE score sees.
\item \textbf{The collective mode is a symmetry direction.} The norm mode is translation along the kink geodesic, and
\textsc{ReLU} commutes with it. For the input's own radius this gives an exact factorization, $\E h_L=\E|x|\cdot\E_{\hat x\sim\text{unif}}h_L(\hat x)$.
Any memory carried by the collective mode is a mixture over this translation. It can therefore be integrated out in one dimension
rather than expanded in cumulants, to the extent that the radial factor decouples from the angular one.
\item \textbf{Two clocks.} The forward state localizes and the backward state delocalizes as $(\tau_l/\tau_L)^3$. The Euler split shows where
the upper-tier (mean--mean) term is reliable and where the mean is a pure border functional.
\end{enumerate}

\subsection{Predictions}
\begin{enumerate}[leftmargin=*]
\item For a fixed He network and $l\ll\sqrt n$, the empirical field law of layer $l+1$ is $N(0,t_l)$ with $t_l\to\cos\theta_l/(1-\cos\theta_l)$. The
unresolved fraction is $\Pr(|r|<1)=\operatorname{erf}(1/\sqrt{2t_l})$, which is $0.227$ at $l=16$ (measured $0.226$).
\item The across-input variance of $\log|h_l|^2$ saturates: $nV_l\to44.6$, not $2+5l$.
\item The internal clock saturates at $t\approx4/V$ once $\tau\gtrsim2\sqrt n$, while the angle between inputs keeps decreasing.
\item The marginal skewness of deep pre-activations is dominated by the collective mode, $\gamma_1\approx\frac32\eta V_{l-1}r$, and the excess
kurtosis is flat in $r$ at $\approx3\eta' V_{l-1}$, rather than following the one-body birth law, which contributes $O(n^{-1})$. Measured at
width $1024$: $\eta=0.71$--$0.72$ at three depths, $\eta'=0.57$--$0.63$, and $R^2\approx0.9$ for the linear law. Confirmed in form; the
efficiency $\eta$ is not yet derived.
\item The backward clock follows $\prod(1-\theta_{k-1}/\pi)$, with exponent $3$ in Fatou time.
\end{enumerate}

\subsection{Open problems}
\begin{enumerate}[leftmargin=*]
\item \textbf{Rates.} Prove Theorem~\ref{thm:selfloc}(2) with explicit rates at fixed $n$, including the compression by the collective mode, and
identify the law of the quenched clock fluctuations.
\item \textbf{The bent-crossing pair term at width $n$.} Compute its size and sign as a function of $\tau$. It is the only genuinely
two-body memory at third order, and it is the term stage 9's hub pairing did not contain.
\item \textbf{Radial--angular factorization beyond the input.} The input radius decouples exactly. How much of $V_l$ (and hence of the
collective memory) decouples at depth, and can the closure be run on the angular part with the radial mixture integrated exactly?
\item \textbf{A modular bound.} Is the closure error controlled by a modular quantity, for example the relative entropy of the layer
state to its quasi-free projection (the negentropy of $z_l$), through a transport inequality on the chain strata?
\item \textbf{The radial efficiency.} Derive $\eta\approx0.71$, the fraction of the across-input log-norm variance that acts as a common
scale on the bulk noise, and its kurtosis analogue $\eta'\approx0.6$. Both are constant across depths $8$--$16$ at width $1024$.
\item \textbf{Beyond the parabolic regime.} For $\tau\gg2\sqrt n$ the radial temperature dominates. Does the two-tier dynamics there
become the hyperbolic (Gauss-type) renormalization conjectured in stage 11?
\end{enumerate}

\section*{Verification}
The checks are run by \texttt{scripts/verify\_stage12.py}; its output is saved next to this note. The network checks use one He network
($n=1024$, $L=16$, seed 5) and $2\times10^5$ standard inputs.
\begin{center}\small
\begin{tabular}{>{\raggedright\arraybackslash}p{5.2cm}>{\raggedright\arraybackslash}p{9.6cm}}\toprule
identity or prediction & result\\\midrule
1. Self-localization (Thm~\ref{thm:selfloc}) & $q_l$, $t_l$, gate variance, $V_l$: see the table in Section~\ref{sec:clock}; gate variance within $0.006$ of $\theta_{l-1}/2\pi$ at all
layers; $\Pr(|r|<1)=0.226$ against $0.227$ at $l=16$\\
2. Field lemma, clock recursion & $\E\Phi(r)\Phi(-r)=\arccos(\frac t{1+t})/2\pi$ to $10^{-8}$ at $t=0.5,2,13$; $\E m_1^2/\E s_1^2$ equals the angle map to $10^{-6}$\\
3. Temperature split (Prop.~\ref{prop:Tsplit}) & $l=16$: $1-q=0.0779$ against $1-e^{-V/4}\cos\theta_{16}=0.0779$; $l=4,8,12$ off by $+0.024$, $+0.010$, $-0.003$
(quenched scatter)\\
4. Fisher--Rao distance (Prop.~\ref{prop:hyper}) & numerical minimization equals $\sqrt2\operatorname{asinh}(|r|/\sqrt2)$ to $10^{-6}$ at $r=0.5,4.29,-6.67$\\
5. Saturation of $V$ (Prop.~\ref{prop:diffusion}) & $(\frac32-\frac{J_2}{2\pi})/\theta^2=0.99671,0.99997$ at $\theta=0.1,0.01$; $nV_l=30.76$, $41.42$, $44.60$ at $l=16,10^2,10^5$; measured
$V_l$ at most $7\%$ above prediction at every layer\\
6. Collective mode (Prop.~\ref{prop:collective}) & trace share $0.008$, $0.026$, $0.058$, $0.093$ against $tV/4=0.005$, $0.022$, $0.053$, $0.086$ at $l=4,8,12,16$\\
7. Collective-mode cumulant law & $\gamma_1=0.0252r$, $0.0302r$, $0.0336r$ at $l=8,12,16$ ($R^2=0.86$, $0.90$, $0.91$; intercepts $\le10^{-4}$) against $\frac32V_{l-1}r$: ratio $\eta=0.72$, $0.72$, $0.71$.
$\gamma_2$ flat in $r$ ($R^2\le0.11$) at $0.040$, $0.053$, $0.059$ against $3V_{l-1}$ (ratio $0.57$--$0.63$). Compression with $\eta$: $\Var r=3.64$, $7.11$, $10.09$ against
$3.79$, $7.07$, $10.14$ measured
\\
8. Ising couplings (Prop.~\ref{prop:ising}) & log odds ratio $0.05614/0.05605$, $0.03382/0.03388$, $0.13388/0.13206$ (exact/first order)\\
9. Wall-density law (Thm~\ref{thm:walldensity}) & $\int\Phi(r)\Phi(-r)dr=1/\sqrt\pi$ to $10^{-10}$; $\E\varphi=\sin(\Theta/2)/\sqrt\pi$ to $10^{-8}$; gate variance $0.022415$ against
$0.022396$ at $t=100$\\
10. Hermite cancellation (Prop.~\ref{prop:hermite}) & $\E\He_2\varphi=-(1+t)^{-3/2}/\sqrt{2\pi}$ to four digits at $t=1,13,100$\\
11. Marginal memory (Section~\ref{sec:gates}) & Gram--Charlier wall terms against the empirical closure error: correlation $0.983$, $0.993$, $0.996$ at $l=8,12,16$;
residual $\le1.7\times10^{-4}$ per $|r|$-bin; error $<5\times10^{-5}$ for $|r|\ge3$\\
12. Stratified third order (Thm~\ref{thm:strata}) & softplus $\beta=5$: $0.1242$ vs $0.1224$ (difference $0.0018\pm0.0028$); $\beta=20$: $0.1677$ vs $0.1698$ ($-0.0022\pm0.0039$)\\
13. Birth law (Prop.~\ref{prop:birth}) & $c_3=0.326458$; $\int k_3^2=0.27002$; $\E k_3^2=0.02904$ against $0.02988$ at $t=13$, $0.01512$ against $0.01523$ at $t=50$\\
14. Backward clock (Prop.~\ref{prop:backclock}) & measured $0.0108$, $0.0800$, $0.2489$, $0.5571$, $0.8798$ against $0.0133$, $0.0941$, $0.2663$, $0.5589$, $0.8743$ at $l=0,4,8,12,15$; power law: $0.00131/0.00122$
($l=10^2..10^3$), $0.00104/0.00102$ ($l=10^3..10^4$)\\\bottomrule
\end{tabular}
\end{center}
The folding angles and clocks: $\theta_l=1.247$, $1.054$, $0.822$, $0.584$, $0.458$, $0.378$ and $t_l^\infty=0.47$, $0.98$, $2.13$, $5.04$, $8.72$, $13.18$ at $l=1,2,4,8,12,16$;
$\tau_{16}=\Fa(\theta_{16})=23.75$; $\tau_*=2.0\sqrt n=64$ at $n=1024$.

\begin{thebibliography}{99}\small
\bibitem{stage11} Stage 11 note: \emph{tilings, Bratteli diagrams and singular foliations} (this repository, \texttt{notes/stage11}).
\bibitem{stage10} Stage 10 note: \emph{the selected Dirac operator} (\texttt{notes/stage10}).
\bibitem{stage9} Stage 9 note: \emph{the mean of a random ReLU network as a central localization} (\texttt{notes/stage9}).
\bibitem{Eldan} R.~Eldan, Thin shell implies spectral gap up to polylog via a stochastic localization scheme, GAFA 23 (2013) 532--569.
\bibitem{Montanari} A.~Montanari, Sampling, diffusions, and stochastic localization, arXiv:2305.10690.
\bibitem{CS} Y.~Cho, L.~K.~Saul, Kernel methods for deep learning, NeurIPS 2009.
\bibitem{HN} B.~Hanin, M.~Nica, Products of many large random matrices and gradients in deep neural networks, Comm.~Math.~Phys.~376 (2020) 287--322.
\bibitem{Szekeres} G.~Szekeres, Regular iteration of real and complex functions, Acta Math.~100 (1958) 203--258.
\bibitem{Milnor} J.~Milnor, \emph{Dynamics in One Complex Variable}, 3rd ed., Princeton 2006 (\S10).
\bibitem{Sudakov} V.~N.~Sudakov, Typical distributions of linear functionals in finite-dimensional spaces of high dimension, Soviet Math.~Dokl.~19 (1978) 1578--1582;
P.~Diaconis, D.~Freedman, Asymptotics of graphical projection pursuit, Ann.~Statist.~12 (1984) 793--815.
\bibitem{Amari} S.~Amari, H.~Nagaoka, \emph{Methods of Information Geometry}, AMS 2000; L.~T.~Skovgaard, A Riemannian geometry of the multivariate normal model, Scand.~J.~Statist.~11 (1984) 211--223.
\bibitem{EP} L.~J.~Emrich, M.~R.~Piedmonte, A method for generating high-dimensional multivariate binary variates, Amer.~Statist.~45 (1991) 302--304;
J.~H.~Macke, P.~Berens, A.~S.~Ecker, A.~S.~Tolias, M.~Bethge, Generating spike trains with specified correlation coefficients, Neural Comput.~21 (2009) 397--423.
\bibitem{Mattis} D.~C.~Mattis, Solvable spin systems with random interactions, Phys.~Lett.~A 56 (1976) 421--422.
\bibitem{Kikuchi} R.~Kikuchi, A theory of cooperative phenomena, Phys.~Rev.~81 (1951) 988--1003.
\bibitem{McC} P.~McCullagh, \emph{Tensor Methods in Statistics}, Chapman and Hall 1987 (Edgeworth series in tensor notation).
\bibitem{LP} A.~Lytova, L.~Pastur, Central limit theorem for linear eigenvalue statistics of random matrices with independent entries, Ann.~Probab.~37 (2009) 1778--1840 (the cumulant expansion).
\bibitem{Stein} C.~Stein, Estimation of the mean of a multivariate normal distribution, Ann.~Statist.~9 (1981) 1135--1151.
\bibitem{Renault} J.~Renault, \emph{A Groupoid Approach to C$^*$-Algebras}, LNM 793, Springer 1980.
\bibitem{FM} J.~Feldman, C.~C.~Moore, Ergodic equivalence relations, cohomology, and von Neumann algebras I, II, Trans.~AMS 234 (1977) 289--359.
\bibitem{BR} O.~Bratteli, D.~W.~Robinson, \emph{Operator Algebras and Quantum Statistical Mechanics 2}, 2nd ed., Springer 1997 (KMS states, quasi-free states).
\bibitem{Yang} G.~Yang, Tensor programs II: neural tangent kernel for any architecture, arXiv:2006.14548 (gradient independence).
\bibitem{PT} I.~F.~Putnam, R.~Trevi\~no, Bratteli diagrams, translation flows and their $C^*$-algebras, arXiv:2205.01537.
\end{thebibliography}
\end{document}
